\chapter{Conformal Manifolds and Exactly Marginal Deformations}
\label{ch:volume-iii-conformal-manifolds-exactly-marginal-deformations}

\section{Families of Conformal Data}

Fix the spacetime dimension \(D\).  A conformal field theory in this volume is
specified by its Hilbert space on \(S^{D-1}\), the action of the conformal
group, the spectrum of primary operators, and the OPE coefficients satisfying
reflection positivity, Ward identities, and crossing.  Denote this complete
data by
\[
  \mathfrak D.
\]
A family of conformal theories parametrized by a smooth manifold
\(\mathcal M_{\rm conf}\) is an assignment
\[
  p\in\mathcal M_{\rm conf}
  \quad\longmapsto\quad
  \mathfrak D(p),
\]
such that correlation functions vary smoothly with \(p\) after a smooth
choice of local operator basis has been made on each coordinate patch.

Changing the basis of local operators changes the coordinate description of
\(\mathfrak D(p)\).  If \(\mathcal O_i(p)\) is a local basis of primaries and
\[
  \mathcal O_i(p)\mapsto U_i{}^j(p)\mathcal O_j(p)
\]
with \(U(p)\) invertible inside a degenerate subspace of quantum numbers, then
the two-point matrix and OPE coefficients transform covariantly.  The point of
\(\mathcal M_{\rm conf}\) is the equivalence class of the complete data under
such basis changes.  This basis dependence is the ordinary coordinate
dependence of a vector bundle of operator spaces over \(\mathcal M_{\rm conf}\).

\begin{definition}[Conformal manifold]
  A conformal manifold is a smooth family \(\mathcal M_{\rm conf}\) of
  conformal data for which every point \(p\in\mathcal M_{\rm conf}\) defines a
  conformal field theory with the same spacetime dimension and the same chosen
  superselection structure, while local operator spaces, scaling dimensions,
  and OPE coefficients vary smoothly in local trivializations.
\end{definition}

At a regular point \(p\), the tangent space \(T_p\mathcal M_{\rm conf}\) is
represented by integrated local deformations generated by scalar primary
operators of dimension \(D\).  Let
\[
  \{\mathcal O_a(x)\}_{a=1}^{m}
\]
be a basis of scalar primaries with \(\Delta_a=D\) at \(p\).  A tangent vector
\[
  v=v^a\partial_a\in T_p\mathcal M_{\rm conf}
\]
is represented infinitesimally by
\begin{equation}
  \delta_v S
  =
  \sum_{a=1}^{m} v^a\int \dd^Dx\,\mathcal O_a(x),
  \label{eq:infinitesimal-marginal-deformation}
\end{equation}
with the understanding that the integrated insertion is defined by a local
renormalization prescription.  Directions generated by descendants or by
field redefinitions are redundant and are quotiented when identifying
\eqref{eq:infinitesimal-marginal-deformation} with a tangent vector.

\begin{figure}[H]
  \centering
  \resizebox{0.92\textwidth}{!}{%
  \begin{tikzpicture}[x=1cm,y=1cm,every node/.style={font=\scriptsize}]
    \tikzset{
      arrow/.style={-{Latex[length=2mm]},line width=0.85pt},
      flow/.style={-{Latex[length=2mm]},line width=0.8pt,orange!85!black},
      locus/.style={blue!65!black,line width=1.1pt},
      op/.style={draw,rounded corners=3pt,fill=blue!5,inner sep=3pt}
    }
    \draw[locus] (-3.8,-0.20)
      .. controls (-2.35,0.90) and (-0.95,0.92) .. (0.15,0.10)
      .. controls (1.35,-0.78) and (2.65,-0.48) .. (3.70,0.38);
    \node[blue!65!black] at (2.00,1.08) {\(\mathcal M_{\rm conf}\)};

    \coordinate (p) at (-0.55,0.47);
    \fill[violet!80!black] (p) circle (2.2pt);
    \node[violet!80!black,above] at (p) {\(p\)};
    \draw[arrow,green!45!black] (p)--++(1.15,-0.55)
      node[midway,below right] {\(v^a\partial_a\)};
    \node[op] at (-1.30,-1.10)
      {\(\displaystyle v^a\int\dd^Dx\,\mathcal O_a\)};
    \draw[arrow] (-1.06,-0.78)--(-0.70,0.18);

    \foreach \x/\y/\dx/\dy in {-3.1/1.0/0.45/-0.22,-2.25/-1.0/0.42/0.25,
      0.75/1.03/0.45/-0.25,2.50/-0.72/0.50/0.18,3.25/1.02/0.35/-0.28} {
      \draw[flow] (\x,\y)--++(\dx,\dy);
    }
    \node[orange!85!black,align=center] at (-3.05,-1.78)
      {beta-function vector field\\vanishes on the conformal locus};
  \end{tikzpicture}
  }
  \caption{At a regular point of a conformal manifold, a tangent vector is
  represented by an integrated dimension-\(D\) scalar primary modulo redundant
  deformations.}
  \label{fig:conformal-manifold-tangent-space}
\end{figure}

\section{Beta Functions and Obstructions}

Let \(\lambda^a\) be local coordinates on the space of marginal couplings near
\(p\).  The deformed action is formally
\begin{equation}
  S_\lambda
  =
  S_p+\sum_a\lambda^a\int\dd^Dx\,\mathcal O_a(x)
  \label{eq:local-marginal-family-action}
\end{equation}
before counterterms are chosen.  A renormalization scheme assigns beta
functions
\[
  \beta^a(\lambda)=\mu\frac{\dd\lambda^a}{\dd\mu},
\]
where \(\mu\) is a reference energy scale.  The conformal locus in this
coordinate patch is the zero set
\[
  \beta^a(\lambda)=0
  \qquad
  \text{for every }a,
\]
with redundant directions divided out.  If this zero set is smooth and the
rank of the obstruction equations is locally constant, it is a conformal
manifold.

The first obstruction to exact marginality is visible in conformal
perturbation theory.  The OPE of two marginal operators has the form
\begin{equation}
  \mathcal O_a(x)\mathcal O_b(0)
  =
  \frac{G_{ab}}{|x|^{2D}}\mathbf 1
  +
  \frac{C_{ab}{}^c}{|x|^D}\mathcal O_c(0)
  +\sum_{\mathcal X}
    \frac{C_{ab}{}^{\mathcal X}}
         {|x|^{2D-\Delta_{\mathcal X}}}\mathcal X(0)
  +\cdots .
  \label{eq:marginal-operator-ope}
\end{equation}
Here \(G_{ab}\) is the two-point matrix, \(C_{ab}{}^c\) are OPE coefficients
projected onto marginal scalar primaries, and the remaining sum runs over
other primaries \(\mathcal X\).  The term proportional to
\(|x|^{-D}\mathcal O_c(0)\) produces a logarithmic divergence when the
integrated pair of insertions approaches coincidence:
\[
  \int_{\epsilon<|x|<L}\dd^Dx\,\frac1{|x|^D}
  =
  {\rm Vol}(S^{D-1})\log\frac L\epsilon .
\]
Consequently, in a scheme adapted to the operator basis,
\begin{equation}
  \beta^c(\lambda)
  =
  -\frac12{\rm Vol}(S^{D-1})\,C_{ab}{}^c\lambda^a\lambda^b
  +O(\lambda^3)
  \label{eq:conformal-perturbation-beta-quadratic}
\end{equation}
up to a convention-dependent normalization of the couplings and operators.
Equation \eqref{eq:conformal-perturbation-beta-quadratic} expresses the
second-order obstruction in terms of OPE data of the undeformed theory.

Higher orders are governed by integrated higher-point functions with
subtractions at coincident points.  Exact marginality means that the
counterterms and coupling redefinitions can be chosen so that the beta
function vector field vanishes along the proposed family and the conformal
Ward identities remain valid at every point.

\section{Metric and Operator Bundles}

At a reflection-positive point, the two-point function of marginal primaries
defines a positive bilinear form:
\begin{equation}
  \left\langle\mathcal O_a(x)\mathcal O_b(0)\right\rangle_p
  =
  \frac{G_{ab}(p)}{|x|^{2D}}.
  \label{eq:zamolodchikov-metric-definition}
\end{equation}
The matrix \(G_{ab}(p)\) is the Zamolodchikov metric on
\(T_p\mathcal M_{\rm conf}\) after redundant directions have been removed.
Contact terms affect local formulas for connections on operator bundles, but
they leave the separated-point coefficient in
\eqref{eq:zamolodchikov-metric-definition} well defined once the operators
are normalized.

As \(p\) varies, primaries with equal spin, charges, and dimension may mix.
Let \(\mathcal V_{\alpha}\to\mathcal M_{\rm conf}\) be the vector bundle
whose fiber at \(p\) is the finite-dimensional space of primaries with a
chosen set of quantum numbers \(\alpha\), assuming no collision with a
different quantum-number sector in the chosen patch.  A connection
\(\nabla_a\) on \(\mathcal V_\alpha\) is defined by subtracting contact
ambiguities from the integrated three-point function
\[
  \int\dd^Dx\,
  \left\langle
    \mathcal O_a(x)\mathcal X_i(y)\mathcal X_j(0)
  \right\rangle .
\]
This connection records the fact that an operator label at one point of
\(\mathcal M_{\rm conf}\) has no absolute meaning until a transport
prescription has been chosen.

The derivatives of scaling dimensions are also CFT data.  For a scalar
primary \(\mathcal X\) with normalized two-point function, the first variation
of its scaling dimension has the schematic form
\begin{equation}
  \partial_a\Delta_{\mathcal X}
  =
  -{\rm Vol}(S^{D-1})\,C_{a\mathcal X\mathcal X}
  \label{eq:dimension-variation-marginal-direction}
\end{equation}
when \(\mathcal X\) is isolated.  In a degenerate subspace,
\eqref{eq:dimension-variation-marginal-direction} becomes a matrix whose
eigenvalues give the first-order splitting.  This is conformal perturbation
theory expressed intrinsically in the operator algebra.

\section{Supersymmetric Protection}

Supersymmetry can impose algebraic conditions that make the obstruction
equations soluble.  In a four-dimensional \(\mathcal N=1\) theory, an
F-term deformation by a chiral primary \(\mathcal W\) has the form
\[
  \delta S
  =
  \lambda\int\dd^4x\,\dd^2\theta\,\mathcal W
  +\bar\lambda\int\dd^4x\,\dd^2\bar\theta\,\overline{\mathcal W}.
\]
It is marginal when the chiral primary has \(R\)-charge \(2\), equivalently
dimension \(3\).  The D-term completion is determined by supersymmetry.
The beta functions are constrained by the superconformal \(R\)-symmetry and by
possible mixing with flavor symmetries.

In four-dimensional \(\mathcal N=4\) super-Yang--Mills, the complex coupling
\[
  \tau_{\rm YM}=\frac{\theta}{2\pi}+\frac{4\pi\ii}{g_{\rm YM}^2}
\]
is an exactly marginal parameter.  The corresponding marginal operator is a
descendant in the stress-tensor multiplet, and its conjugate supplies the
anti-holomorphic tangent direction.  The local metric on this conformal
manifold is fixed by supersymmetry up to the normalization of the stress
tensor multiplet.  For gauge algebra \(\mathfrak{su}(N)\), the local upper
half-plane coordinate \(\tau_{\rm YM}\) will be further identified by
duality transformations discussed in the next chapter.

\begin{figure}[H]
  \centering
  \resizebox{0.88\textwidth}{!}{%
  \begin{tikzpicture}[x=1cm,y=1cm,every node/.style={font=\scriptsize}]
    \tikzset{
      arrow/.style={-{Latex[length=2mm]},line width=0.85pt},
      box/.style={draw,rounded corners=4pt,line width=0.8pt,fill=blue!4,
        minimum width=2.9cm,minimum height=0.82cm,align=center}
    }
    \node[box] (marginal) at (-4.2,0.85)
      {dimension-\(D\)\\scalar primary};
    \node[box] (integrated) at (-4.2,-0.85)
      {integrated\\local insertion};
    \node[box] (ope) at (0,0.85)
      {OPE coefficients\\\(C_{ab}{}^c\)};
    \node[box] (beta) at (0,-0.85)
      {beta functions\\\(\beta^a(\lambda)\)};
    \node[box] (locus) at (4.2,0)
      {conformal locus\\\(\beta=0\)};
    \draw[arrow] (marginal)--(integrated);
    \draw[arrow] (marginal)--(ope);
    \draw[arrow] (ope)--(beta);
    \draw[arrow] (integrated)--(beta);
    \draw[arrow] (beta)--(locus);
    \node[green!45!black,align=center] at (0,-2.0)
      {exact marginality is the vanishing of all obstruction components};
  \end{tikzpicture}
  }
  \caption{Marginal operators become coordinates on a conformal manifold
  precisely when the beta-function obstructions vanish after quotienting
  redundant directions.}
  \label{fig:exactly-marginal-obstruction-flow}
\end{figure}
